\section{Fouille d'arguments et extraction de structures par LLM}
Lawrence et Reed définissent le champ autour de l'identification de structures d'inférence exprimées en langage naturel \cite{ref05}. Stab et Gurevych combinent étiquetage séquentiel et optimisation linéaire entière des composantes et relations \cite{ref07}. Lauscher et collègues ajoutent des annotations argumentatives au corpus scientifique Dr. Inventor \cite{ref08}. Ces approches rendent explicite la distinction entre unité textuelle, fonction argumentative et relation.

La revue de Li et collègues décrit l'évolution des tâches et des évaluations avec les LLM, ainsi que des difficultés de contexte long et de fiabilité des juges \cite{ref06}. Pour notre prototype, cela justifie de séparer proposition initiale, contrôle d'ancrage et critique sémantique. Faire juger une sortie par un autre appel du même modèle peut détecter des incohérences ; cela ne garantit pas une indépendance des erreurs.

GraphRAG construit un graphe d'entités puis des résumés de communautés pour répondre à des questions globales \cite{ref13}. LightRAG associe structures graphiques et recherche à deux niveaux \cite{ref14}. Nous utilisons ces travaux comme contraste : l'association entre deux concepts ne précise pas automatiquement lequel soutient l'autre. Notre schéma exige un type de relation et une provenance textuelle, avec un statut distinct pour les éléments insuffisamment étayés.
